\section{결론}

이 논문은 TSI 제안을 하나의 formal core로 전환했다. Integrated structural imagination state는 shared tagged carrier, typed relation, simplicial complex, metric, filtration, mass, 그리고 action-conditioned transition을 함께 갖는다. Structural equivalence는 label, simplex, distance, filtration, mass, generator relation을 모두 보존하는 bijection으로 정의되며 equivalence relation임을 증명했다. Partial transition에 대해서는 topology, geometry, relation, filtration, mass의 다섯 survivor layer와 turnover를 분리하고, 선택한 layer의 보존과 유한 composition closure를 증명했다. Quotient-level Lipschitz update에는 명시적인 rollout bound를 제시했다.

이 결과는 TSI가 intelligence에 필요충분하다는 주장이나 cognitive validity를 확립하지 않는다. 대신 각 직관을 state object, 명시적 가정, 증명 가능한 정리, 그리고 실패 경계로 분해한다. 추후 연구를 통해 이 formal contract의 적용 범위와 경험적 유용성을 검증할 수 있다.

다음 결정적 비교는 factorized head와 generic sparse head의 비교가 아니다. Entity의
생성과 삭제, identity 유지 및 relation 재구성 아래에서 동일한 dependency graph를
회수한 두 model을 비교해야 한다. 이 비교가 어떤 효과를 식별하려면 두 arm이 동일한
observable information을 받아야 하며 representational commitment만 달라야 한다.
Identity 또는 typed relation을 한 arm에만 input으로 제공하면 결과는 information
asymmetry의 산물이 되며, 이는 원래 finite-family audit에서 진단한 것과 같은 종류의
confound다. Turnover quantity \(b(p)\), 그 zero criterion과 survivor-only preservation의
vacuity 명제는 이미 여기서 formalize되어 있다. 따라서 turnover에 초점을 둔 연구가 이
설계의 가장 작은 self-contained instance다.
